\subsection*{先给保证指定对象和决定位置}

从最早的本地动作 \code{purchase.create.v1} 开始。它在同一个事务中仲裁 $S$，保存原输入，扣减本地余额，建立实际 P/O/D 与原结果。对于一项有效且已提交的 $S$，同参重试取得原 P/O/D，异参来件得到冲突。两项不同的合法 $S$ 即使商品和金额相同，也各自参与首次仲裁。这些性质依赖同一权威事务中的唯一性、原子性和已确认历史的保存。

例如，S1 成功时余额由 10000 变为 7000，随后 S2 成功，当前余额成为 4000。恢复 S1 时返回原来的 7000，账户查询返回当前 4000。两份结果各自有确定的对象：前者是 S1 的成交快照，后者是查询时的账户状态。因此，接口说明中的重复成功应包含原购买、原输入和历史结果之间的固定关系。

外部动作 \code{purchase.create.external.v2} 把资金执行移出了这个事务域。$A_J$ 确定本系统只为 S1 建立一份 C1，$R_J$ 把同 C1 的可信成功证据连接到原 O1 和 EV1。平台对原身份的仲裁决定多次到达能否合到同一 X1。到了履约域，I 记录 EV1 的原接收结果，B 以 O1 的初次履约为仲裁对象，两者共同指向已经提交的 F501。EV1 再次到达时，消费者据此返回原 F501；另一事件报告同一 O1 时，业务屏障也会找到这份初次履约。

\CFifteenFigure{ch15-01-domains}{同一购买经过独立决定，身份、原结果和继续责任连接这些提交}{fig:ch15-domains}

因此，把整条链称为恰好一次时，应给出具体计数。客户端请求可以发送多次，工作者可以领取多次，平台请求可以到达多次，EV1 可以发布多次。同一个 $S$ 的权威受理绑定至多一份，同 C 的资金成功在下游身份前提内至多一份，同 B 的初次履约至多一份。有效工作最终得到执行机会，且业务判定允许成功时，这些位置才能各自形成一份成功事实。

消费者提交 F501 后、ACK 发出以前退出，就是一条直观历史。消息系统再次投递 EV1，Q 按 I 恢复原接收结果，按 B 找到 F501，随后再次确认消息。投递计数为二，已提交初次履约计数为一。若 F501 之后还要调用承运商，消费者先保存 \code{dispatch:initial:O1} 及原发送内容。承运商在自己的接收事务中以该身份仲裁并保存运单结果，发送者失联后便可查询或重发原身份，取得同一运单。承运商若只接受普通调用，已经发出却未取得回执的那次发送就进入持久核查，由账单、运单查询或人工回执确定是否已经发货。

\CFifteenFigure{ch15-17-counts}{同一 EV1 投递两次，I/B 将两次接收连接到同一 F501}{fig:ch15-counts}

\subsection*{平台能力怎样改变继续方式}

设本地 C1 未决，平台请求曾经发出，调用方只看见超时。第一种平台同时提供稳定身份和可信结果查询。恢复者从 C1 取得完整 \code{provider\_scope}、\code{pay-C1} 与原支付输入，查询到 X1 便核验并保存结果。需要重新发送时，每一个可能生效的到达都落在平台对原身份的保护范围中，才继续使用同一键。

这里的范围包括商户、环境和动作。本书主线的完整范围为
\begin{BookCode}
simulator:merchant-7:test:charge.v1
\end{BookCode}
把测试商户的 \code{pay-C1} 换到另一个商户，或者把 charge 动作的键用于 refund，接收端可能面对另一个命名空间。平台回来的账户、C1、币种、金额和 X1 也要与本地关联逐项匹配，再由可信成功证据生成 $R_J$。

第二种平台只有稳定身份，没有独立查询。平台在有效范围内以同键重试返回原结论，恢复者就可以从重试响应取得证据。假如所有响应永久丢失，或者逐键记录先被回收，系统可能一直无法获知结果。平台效果的上限由原身份保护决定，本地形成终态还依赖响应最终到达并保存。

第三种平台只提供查询。查到可信 X1，可以结束未知。查空时，旧请求却可能继续留在平台队列里。若此时新建另一个付款身份，旧请求和新请求都可能成功。恢复者继续核对原关联，或取得能够排除旧请求后来成功的终结依据，再按业务规则决定新的尝试。

第四种平台既没有可依赖的身份仲裁，也没有足够结果证据。请求发出后失联，恢复者保存可能已执行的事实和原输入，限制自动发送，等待账单或其他可靠证据。业务可以向用户显示待核查，关闭该支付方式的新受理，或者采用事先公布的人工处理。每种选择都把未知工作交给一个持续存在的责任位置。

\begin{table}[htbp]
\centering\small
\begin{tabularx}{\linewidth}{@{}p{.23\linewidth}XX@{}}
\toprule
下游能力 & C1 未决时的继续依据 & 形成终态的进展条件\\
\midrule
原身份保护与可信查询 & 查询原关联，在保护范围内重发原键 & 原结果最终可取得并保存\\
仅原身份保护 & 在适用范围内重发原键 & 原终态响应最终返回\\
仅可信查询 & 核对原关联与账本 & 查空后取得旧路径终结依据\\
两项能力都不足 & 保存未知与核查责任 & 后续取得足够资金证据\\
\bottomrule
\end{tabularx}
\caption{平台能力、恢复动作与终态条件分别说明}\label{tab:ch15-capabilities}
\end{table}

平台文档中的期限也进入这份说明。Stripe 在保留原键记录期间比较请求参数，以识别同键异参。其 API v1 文档允许在键至少保存 24 小时后清理，清理后再次使用该键，会建立新请求~\cite{stripe-v1-idempotency}\footnote{Stripe 文档中的 API v1、API v2 是平台接口分类，与本书本地购买 v1、外部购买 v2 的动作版本分别命名。实际平台的保护期限与原身份关闭规则按其接口分别核对。}。本地为 C1 保存七天记录，发送期限却受平台逐键规则约束。发送者根据实际平台的受理窗口、在途上界或同域关闭能力决定是否继续发送。

\CFifteenFigure{ch15-02-provider}{身份保护使安全重发合流，查询提供结果证据，旧在途请求另有终结依据}{fig:ch15-provider}

\subsection*{故障范围扩大时，哪份历史继续保存}

\begin{table}[htbp]
\centering\small
\caption{故障模型与权威来源的演进}\label{tab:ch15-fault-evolution}
\begin{tabularx}{\linewidth}{@{}p{30mm}X X@{}}
\toprule
变化或故障 & 延续决定的来源 & 入口的具体判断\\
\midrule
单库进程崩溃 & 同一事务与持久恢复日志中的 S、业务变化和原结果 & 从原绑定恢复，未提交事务回滚\\
异步跨网调用 & A 的固定 C/T 与独立平台账本 & 未知保留原身份，核验后共同保存结果\\
副本切换、快照截断 & 已提交前缀及完整同边界快照 & 接班者恢复并应用到有效读取边界\\
结果与身份回收 & 活动身份依据，或不可撤回关闭的会话资格 & 未退休结果回放，关闭旧入口拒绝\\
购买库旧备份回退 & 独立 J 的 $A_J/R_J$、会话和目录历史，平台与消费事实 & 冻结受影响范围，重建原关联与未决台账\\
\bottomrule
\end{tabularx}
\end{table}

最后一行允许 W 的投影变旧，权威 J、会话资格 $\Gamma$ 和目录 D 必须提供跨过此次回退的完整来源。恢复者按原 S/P/C/T 和 O/EV 补写缺失行，完整核对已有内容，平台账本用于核对 X1，履约域用于核对 F501。共同来源也丢失时，受影响入口保持冻结，可靠归档或其他足够证据恢复以后再决定开放范围。

目录授予的路由代次控制后继地域资格和新编号。旧 S1 的恢复继续指向原 C1/\code{pay-C1}；真正的新购买 S2 才取得新身份。受限开放在同一权威序列中核验允许子集、完整受理与会话历史，以及旧外部路径的保护或关闭依据，通过后再登记新会话。

一份对外说明于是可以落到具体句子上：在 E1 的固定恢复期限内，携带原完整 $S$ 和原输入的已受理购买，从原绑定恢复受理信息与已保存的原结果。输入冲突保留原绑定。E1 关闭后，购买提交入口返回过期，授权查询与服务端 C/T 恢复继续沿原关系工作。说明同时列出适用的部署保存范围、平台身份窗口和结果证据条件，调用方据此保存任务、安排查询并识别需要介入的未知。
\FloatBarrier
